% Options for packages loaded elsewhere
\PassOptionsToPackage{unicode}{hyperref}
\PassOptionsToPackage{hyphens}{url}
\documentclass[
]{article}
\usepackage{xcolor}
\usepackage{amsmath,amssymb}
\setcounter{secnumdepth}{-\maxdimen} % remove section numbering
\usepackage{iftex}
\ifPDFTeX
  \usepackage[T1]{fontenc}
  \usepackage[utf8]{inputenc}
  \usepackage{textcomp} % provide euro and other symbols
\else % if luatex or xetex
  \usepackage{unicode-math} % this also loads fontspec
  \defaultfontfeatures{Scale=MatchLowercase}
  \defaultfontfeatures[\rmfamily]{Ligatures=TeX,Scale=1}
\fi
\usepackage{lmodern}
\ifPDFTeX\else
  % xetex/luatex font selection
\fi
% Use upquote if available, for straight quotes in verbatim environments
\IfFileExists{upquote.sty}{\usepackage{upquote}}{}
\IfFileExists{microtype.sty}{% use microtype if available
  \usepackage[]{microtype}
  \UseMicrotypeSet[protrusion]{basicmath} % disable protrusion for tt fonts
}{}
\makeatletter
\@ifundefined{KOMAClassName}{% if non-KOMA class
  \IfFileExists{parskip.sty}{%
    \usepackage{parskip}
  }{% else
    \setlength{\parindent}{0pt}
    \setlength{\parskip}{6pt plus 2pt minus 1pt}}
}{% if KOMA class
  \KOMAoptions{parskip=half}}
\makeatother
% definitions for citeproc citations
\NewDocumentCommand\citeproctext{}{}
\NewDocumentCommand\citeproc{mm}{%
  \begingroup\def\citeproctext{#2}\cite{#1}\endgroup}
\makeatletter
 % allow citations to break across lines
 \let\@cite@ofmt\@firstofone
 % avoid brackets around text for \cite:
 \def\@biblabel#1{}
 \def\@cite#1#2{{#1\if@tempswa , #2\fi}}
\makeatother
\newlength{\cslhangindent}
\setlength{\cslhangindent}{1.5em}
\newlength{\csllabelwidth}
\setlength{\csllabelwidth}{3em}
\newenvironment{CSLReferences}[2] % #1 hanging-indent, #2 entry-spacing
 {\begin{list}{}{%
  \setlength{\itemindent}{0pt}
  \setlength{\leftmargin}{0pt}
  \setlength{\parsep}{0pt}
  % turn on hanging indent if param 1 is 1
  \ifodd #1
   \setlength{\leftmargin}{\cslhangindent}
   \setlength{\itemindent}{-1\cslhangindent}
  \fi
  % set entry spacing
  \setlength{\itemsep}{#2\baselineskip}}}
 {\end{list}}
\usepackage{calc}
\newcommand{\CSLBlock}[1]{\hfill\break\parbox[t]{\linewidth}{\strut\ignorespaces#1\strut}}
\newcommand{\CSLLeftMargin}[1]{\parbox[t]{\csllabelwidth}{\strut#1\strut}}
\newcommand{\CSLRightInline}[1]{\parbox[t]{\linewidth - \csllabelwidth}{\strut#1\strut}}
\newcommand{\CSLIndent}[1]{\hspace{\cslhangindent}#1}
\setlength{\emergencystretch}{3em} % prevent overfull lines
\providecommand{\tightlist}{%
  \setlength{\itemsep}{0pt}\setlength{\parskip}{0pt}}
\usepackage{graphicx}
\usepackage{float}
\graphicspath{{figures/out/}}
\usepackage{bookmark}
\IfFileExists{xurl.sty}{\usepackage{xurl}}{} % add URL line breaks if available
\urlstyle{same}
\hypersetup{
  hidelinks,
  pdfcreator={LaTeX via pandoc}}

\author{}
\date{}

\begin{document}

\section{Problem-Shaped Virtual Machines for Exact Local Transformer
Execution}\label{problem-shaped-virtual-machines-for-exact-local-transformer-execution}

\textbf{Aviraj Khare}\\
Email: \texttt{avirajkhare00@gmail.com}

\subsection{Abstract}\label{abstract}

This paper argues that small local models are most useful on exact tasks
when they operate inside a task-specific exact machine instead of
predicting final answers or imitating a broad virtual machine. I call
that machine a \textbf{problem-shaped virtual machine (PSVM)}. The
repository demonstrates the pattern across browser-local Sudoku, receipt
and invoice total extraction, Tally-style voucher extraction, and a
smaller Weiqi prototype (Khare 2026). In every case, code remains
authoritative for legality, rollback, arithmetic, schema checks, and
final emission; the model only ranks legal branches or candidates. The
contribution is therefore systems-level: a concrete artifact for
browser-local, verifier-backed AI. The paper does not claim that PSVMs
already beat strong baselines or scale to arbitrary program execution.

\subsection{Core Claim}\label{core-claim}

For narrow exact tasks, the right learned target is usually not the
final answer and not a full machine trace. It is the ranking problem
induced by the smallest exact machine that can still run the task
correctly.

That claim sits between two common extremes. One-shot answer prediction
hides the exact intermediate decisions, making supervision weak and
failure analysis vague. Full-machine emulation forces the model to spend
capacity on generic machine bookkeeping that the task does not need.
PSVMs aim at the middle: keep exact semantics in code, but expose a
small legal frontier on which the model can actually help.

This is a repository-backed systems paper, not a finished benchmark
study. The point is to make the pattern concrete in working software.
OR-Tools and PySAT provide exact search and constraint solving (Google
OR-Tools Team 2026; Ignatiev, Morgado, and Marques-Silva 2018), while
Transformers.js and ONNX Runtime make browser-local inference practical
(Boreiko 2026; Microsoft 2026). LayoutParser shows how OCR-centric
document pipelines can be structured (Shen et al. 2021). The gap this
repository targets is the combination: exact task runtimes, local
learned guidance, canonical candidate surfaces or traces, and browser
deployment in one inspectable artifact.

\subsection{What a PSVM Is}\label{what-a-psvm-is}

A PSVM is a small exact machine for one task family. It keeps only the
state and transitions that matter for that task, plus an explicit place
where a model can help. In practice, a PSVM has five parts:

\begin{itemize}
\tightlist
\item
  exact state
\item
  legal next actions or legal candidate set
\item
  exact verifier
\item
  canonical trace
\item
  model hook for ranking ambiguity
\end{itemize}

Across tasks, the execution loop is:

\texttt{task\ instance\ -\textgreater{}\ exact\ state\ -\textgreater{}\ legal\ frontier\ -\textgreater{}\ model\ scores\ ambiguity\ -\textgreater{}\ runtime\ verifies\ and\ applies\ -\textgreater{}\ new\ state}

Let \(S_t\) be the exact state at step \(t\), \(V(S_t)\) the valid next
actions, and \(m_\theta(S_t, a)\) the model score for legal action
\(a\). Then the guided choice is

\[
a_t^* = \arg\max_{a \in V(S_t)} m_\theta(S_t, a)
\]

and the runtime applies the exact transition

\[
S_{t+1} = \delta(S_t, a_t^*)
\]

The separation of responsibility is the point:

\begin{itemize}
\tightlist
\item
  the runtime defines \(V(S_t)\) and \(\delta\)
\item
  the verifier decides whether an action is legal
\item
  the model only ranks choices inside the legal set
\end{itemize}

That pattern applies both to search tasks, where actions may include
branch and rollback, and to extraction tasks, where actions may be
candidate selections. The design rule is simple: trim the machine to the
problem.

\subsection{Evidence from the
Repository}\label{evidence-from-the-repository}

The repository demonstrates the PSVM pattern across several domains
rather than only one demo.

\subsubsection{Sudoku}\label{sudoku}

Sudoku is the cleanest case because it is explicit exact search. The
runtime owns candidate generation, contradictions, rollback, and
backtracking. The model only helps order legal branches. On the current
hard preset \texttt{Forum\ hardest\ 1106\ ·\ r365}, the shipped
regret-trained transformer reduced ranked branch decisions from 7,523 to
5,133, backtracks from 86,376 to 57,057, and wall time from 96.99
seconds to 64.94 seconds relative to the imitation-trained transformer
on the same exact runtime. The rules did not change. Only branch
ordering changed.

That result is suggestive, not sufficient. It is a single-preset
comparison inside the repository, not a baseline study against stronger
exact solvers. But it demonstrates the intended role of learning in a
PSVM: not free-running execution, but value estimation over exact
states.

\subsubsection{Receipt, Invoice, and Tally
Extraction}\label{receipt-invoice-and-tally-extraction}

The document lanes show the same separation outside search. For receipt
and invoice total extraction, the runtime first enumerates legal money
candidates from OCR or structured text. The model ranks those
candidates, and a deterministic emitter produces the selected value. The
model is not asked to invent totals from scratch.

The Tally lane extends the same idea to a larger structured state. The
runtime first classifies voucher family, then builds schema-aligned
field candidates, and finally applies a deterministic resolver with
arithmetic, schema, and GST consistency checks. The learned selector
only ranks candidates inside that legal surface. This is the main
document-extraction lesson of the repository: parser quality, candidate
recall, and exact resolution matter at least as much as model size.

\subsubsection{Browser-Local Deployment}\label{browser-local-deployment}

The artifact is also a browser-local systems demo, not just a training
directory. It ships worker-based inference paths, ONNX exports, exact
runtimes, and interactive demos that make the model/runtime split
visible. That matters because latency, inspectability, and deterministic
fallback behavior are part of the thesis, not afterthoughts.

\subsection{Positioning}\label{positioning}

The repository already supports three defensible claims.

\begin{itemize}
\tightlist
\item
  Browser-local exact runtimes with local learned guidance are practical
  for narrow tasks.
\item
  Canonical legal frontiers can be built for both search and extraction
  tasks.
\item
  Small local models are useful as rankers over those legal frontiers
  while exact code keeps truth.
\end{itemize}

It does \textbf{not} yet justify stronger claims such as:

\begin{itemize}
\tightlist
\item
  PSVMs outperform strong classical or hybrid baselines across tasks.
\item
  the learned model can replace the exact runtime end to end.
\item
  the approach scales from narrow task families to arbitrary program
  execution.
\end{itemize}

These are precisely the places where criticism is fair. The current
artifact is strongest as a systems paper or research note, not as a
finished empirical paper.

\subsection{What Would Make It
Stronger}\label{what-would-make-it-stronger}

Three additions would materially strengthen the paper: compare PSVM
ranking against one-shot answer prediction and broader trace
formulations; evaluate exactness rather than only loss, using solve
rate, illegal action rate, verifier failure rate, search cost, trace
length, and browser latency; and add stronger baselines, especially for
Sudoku and document extraction. If those experiments are strong, the
PSVM claim becomes much easier to defend as more than a design argument.

\subsection{Conclusion}\label{conclusion}

The paper's main point is narrower than the title makes it sound. It is
not that local transformers can already execute arbitrary code exactly.
It is that many narrow exact tasks have a small semantic core, and local
models are more useful when trained on that core than when asked either
to guess final answers or imitate broad machine traces.

That is what a PSVM is: the smallest exact machine that still preserves
the task's truth conditions while exposing the choices that are
genuinely ambiguous.

\subsection{AI Assistance Disclosure}\label{ai-assistance-disclosure}

This paper was drafted and edited with AI assistance. The author
reviewed the implementation references and is responsible for the
claims, wording, and technical judgments.

\protect\phantomsection\label{refs}
\begin{CSLReferences}{1}{0}
\bibitem[\citeproctext]{ref-transformersjs}
Boreiko, Xenova. 2026. {``Transformers.js.''}
\url{https://github.com/huggingface/transformers.js}.

\bibitem[\citeproctext]{ref-ortools}
Google OR-Tools Team. 2026. {``OR-Tools.''}
\url{https://developers.google.com/optimization}.

\bibitem[\citeproctext]{ref-pysat}
Ignatiev, Alexey, António Morgado, and Joao Marques-Silva. 2018.
{``PySAT: A Python Toolkit for Prototyping with SAT Oracles.''}
\emph{SAT 2018}. \url{https://doi.org/10.1007/978-3-319-94144-8_26}.

\bibitem[\citeproctext]{ref-khare2026trl}
Khare, Aviraj. 2026. {``Transformer Runtime Lab.''} Zenodo.
\url{https://doi.org/10.5281/zenodo.19087723}.

\bibitem[\citeproctext]{ref-onnxruntime}
Microsoft. 2026. {``ONNX Runtime.''} \url{https://onnxruntime.ai/}.

\bibitem[\citeproctext]{ref-layoutparser}
Shen, Zejiang, Ruochen Zhang, Melissa Dell, Benjamin Charles Germain
Lee, Jacob Carlson, and Weining Li. 2021. {``LayoutParser: A Unified
Toolkit for Deep Learning Based Document Image Analysis.''}
\emph{International Conference on Document Analysis and Recognition}.
\url{https://doi.org/10.1007/978-3-030-86331-9_7}.

\end{CSLReferences}

\end{document}
